\section*{Introduction}

Symmetry is the organizing principle of machine learning for atomistic systems. Networks whose internal features transform equivariantly under the symmetries of three-dimensional space have become the default architecture for predicting the properties of molecules and crystals, from interatomic potentials\cite{batzner2022nequip,musaelian2023allegro,batatia2022mace} to equivariant transformers\cite{liao2023equiformer,liao2024equiformerv2}, and the community describes them almost interchangeably as ``E(3)-equivariant''.

How much of that symmetry should be hard-wired, and how much left to be learned from data, is one of the field's live disputes\cite{creed2026six}. Non-equivariant architectures trained with augmentation close the data-efficiency gap given enough epochs\cite{brehmer2024scale}; unconstrained potentials trained on large corpora can beat physically constrained ones on accuracy and speed\cite{bigi2026unconstrained}; models that are only approximately rotation-equivariant suffer negligible consequences in realistic bulk simulation\cite{langer2024broken}; and the opposing evidence, that equivariance pays increasingly at scale, is empirical rather than structural\cite{ngo2025scaling}. The dispute has been conducted largely on one axis, rotations, and in one currency, accuracy. There is a precedent for what happens on the other kind of axis. Many recent potentials predict forces directly rather than as energy gradients, trading an exact conservation law for speed on the argument that conservation can be learned; assessed against the physics rather than the test error, those models turn out to break geometry optimization and molecular dynamics\cite{bigi2025forces}.

This paper opens a second such axis. The Euclidean group E(3) contains improper operations: reflections and the spatial inversion $I\colon\, \mathbf{r} \mapsto -\mathbf{r}$. A network can be exactly equivariant to all rotations, SO(3), while carrying no information about parity at all, and a substantial class of deployed models is built this way, including the eSCN lineage\cite{passaro2023escn} and EquiformerV2. The distinction is invisible at the configuration surface: such models expose a maximum angular degree, not a parity flag. It is also not a historical artefact of one architecture. The current generation of universal potentials is built on the same spherical-channel lineage\cite{fu2025esen,wood2025uma}, and the newest member of the transformer family declares equivariance to SE(3) in its title\cite{equiformerv3}, a group that contains no improper operation at all. Whether a model respects improper operations is, in practice, an implementation detail that neither model cards nor benchmarks report. On this axis the currency is not accuracy. A model that is slightly wrong about a rotation makes a slightly wrong prediction. A model that is blind to parity can make predictions that are not approximately wrong but physically impossible.

The physics supplies an unusually clean instrument for making this precise. Neumann's principle\cite{nye1985crystals} requires every macroscopic property tensor of a crystal to be invariant under each symmetry operation of the crystal's point group. The piezoelectric tensor $e_{ijk}$, which couples polarization to strain, is a parity-odd rank-3 tensor: under inversion it changes sign. For a centrosymmetric crystal, one whose space group contains the inversion, invariance and sign reversal are simultaneously required, so the tensor must equal its own negative. The consequence is an algebraic identity, not an empirical trend: the piezoelectric tensor of every centrosymmetric crystal is exactly zero. The same holds for every parity-odd tensor, including, in the electric-dipole approximation, the second-harmonic-generation tensor that drives nonlinear-optical materials discovery\cite{optixnet2026}. Since 92 of the 230 space groups are centrosymmetric, this yields a large, label-free test population on which any nonzero prediction is known, with certainty and without reference data, to be impossible. It converts the abstract question of how much symmetry to build in into a measurement.

Whether an equivariant network can violate the identity is decided by its feature algebra, and the answer is not a matter of degree. A parity-labelled network must predict exactly zero on a centrosymmetric crystal, at any parameter values, random initialization included. Erase the labels and the network remains exactly rotation-equivariant, and loses no accuracy where the constraint is inactive, but nothing in it forbids the impossible output. We prove both statements below.

Three adjacent lines of work make the question timely without answering it. A literature on symmetry breaking starts from Curie's principle\cite{curie1894symmetry}, that an equivariant model cannot produce outputs less symmetric than its inputs, and develops relaxations permitting such outputs when the physics demands them\cite{smidt2021symmetry,xie2024symmetrybreaking,kaba2023symmetry,prob2025symmetry}. We probe the complement: for a parity-odd tensor on a centrosymmetric crystal there is no symmetry breaking to model, the answer is exactly zero, and a model that escapes the constraint has not gained flexibility, it has made an error. A second line proves that an E(3)-\emph{invariant} model cannot distinguish a molecule from its mirror image, and so cannot represent chirality\cite{dumitrescu2024chirality}. That is the opposite failure of parity bookkeeping, and a claim about what a model cannot see rather than about what it can wrongly emit. Third, a growing family of crystal-tensor predictors enforces the correct symmetry at the output, through space-group-informed modules, canonicalization, or symmetry-adapted decompositions\cite{yan2024gmtnet,goectp2024,lu2025eatgnn,piezotensornet2024,ceitnet2026}. These are readout-level repairs, and whether the features beneath them can represent the constraint at all is a prior question, one that becomes urgent as property heads are attached to the embeddings of pretrained universal potentials\cite{hung2024dielectric}, where the backbone's symmetry group is fixed long before any downstream target is chosen. Symmetry has also begun to be used to build benchmarks for atomistic models\cite{parackal2026symmetry}; our test is of that kind, but needs no reference calculation.

We test this on 2,000 spglib-verified centrosymmetric insulators, a population on which the true piezoelectric tensor is zero by symmetry, so that any nonzero prediction is a certified error obtained without labels and without a reference calculation. A short group-theoretic account, given in full below, fixes what each symmetry class can guarantee and predicts in advance, with no free parameters, which of these crystals a rotation-only model is nonetheless forced to get right. The measurement is a matched-pair ablation. NequIP, Allegro and MACE are each built in two arms that share architecture, hyperparameters, data, splits and seeds and differ only in the parity labelling of their features, and unmodified EquiformerV2 joins them as a deployed rotation-only representative. The parity-labelled arms stay below the detection threshold on every crystal, at a median of $3\times10^{-7}$\,C\,m$^{-2}$; the rotation-only arms predict piezoelectric responses of ordinary magnitude for roughly nine crystals in ten, landing within one to two percent of the predicted ceiling. Retraining on the exact answer, up-weighting it a hundredfold and symmetrizing at test time all fail to remove the violations, which is a negative result for the proposition that symmetry can be learned from data in the regime where the symmetry is exact. Finally, an audit of fifteen released architectures finds that only six type parity in their features. Three of the remaining nine are rotation-equivariant, and their features carry no parity information at all, so the guarantee cannot be recovered from them; it can only be imposed from outside, one quantity at a time, using symmetry knowledge the model does not itself have. The obstruction therefore belongs to a symmetry class that is in current use, and not to any single model.
